% Self-contained working report. Compile from the repository root: pdflatex -output-directory=docs/report docs/report/Report.tex (twice).
% Refresh/check numerical macros with docs/research/build_report_evidence.py.
\documentclass[11pt,a4paper]{article}
\usepackage[T1]{fontenc}
\usepackage[utf8]{inputenc}
\usepackage[margin=2.1cm]{geometry}
\usepackage{array,booktabs,longtable,tabularx}
\usepackage{amsmath}
\usepackage{xcolor}
\usepackage{enumitem}
\usepackage{xurl}
\usepackage{hyperref}
\definecolor{Page}{HTML}{FFFFFF}
\definecolor{Ink}{HTML}{000000}
\definecolor{Link}{HTML}{000000}
\definecolor{Rule}{HTML}{000000}
\pagecolor{Page}\color{Ink}
\hypersetup{colorlinks=true,urlcolor=Link,linkcolor=Link,citecolor=Link,
  pdftitle={Few-shot OCT: results, open gaps, and research plan},pdfauthor={Lam}}
\setlength{\parindent}{0pt}
\setlength{\parskip}{0.48em}
\setlength{\emergencystretch}{2em}
\setlist{leftmargin=*,itemsep=0.2em,topsep=0.2em}
\renewcommand{\arraystretch}{1.2}
\newcolumntype{Y}{>{\raggedright\arraybackslash}X}
\newcommand{\file}[1]{\nolinkurl{#1}}
\newcommand{\ev}[1]{\textcolor{Link}{\textsf{[#1]}}}
\newcommand{\pp}{percentage points}

% BEGIN VERIFIED NUMBERS (build_report_evidence.py)
\newcommand{\BasicAcc}{88.85}
\newcommand{\BasicCi}{86.75--90.83}
\newcommand{\BasicCommon}{890}
\newcommand{\BasicFone}{0.864}
\newcommand{\BasicKappa}{0.853}
\newcommand{\BasicMacro}{84.56}
\newcommand{\BasicMcc}{0.861}
\newcommand{\BasicRare}{18}
\newcommand{\BasicRareMacro}{81.0}
\newcommand{\BasicRci}{0.760}
\newcommand{\BasicSeedMean}{92.93}
\newcommand{\BasicSeedSd}{1.85}
\newcommand{\GanAcc}{92.95}
\newcommand{\GanCi}{91.19--94.62}
\newcommand{\GanCommon}{931}
\newcommand{\GanFone}{0.901}
\newcommand{\GanKappa}{0.907}
\newcommand{\GanMacro}{89.16}
\newcommand{\GanMcc}{0.909}
\newcommand{\GanRare}{19}
\newcommand{\GanRareMacro}{86.0}
\newcommand{\GanRci}{0.814}
\newcommand{\GanSeedMean}{92.59}
\newcommand{\GanSeedSd}{1.29}
\newcommand{\HeadMinutes}{34.60}
\newcommand{\NoneAcc}{90.31}
\newcommand{\NoneCi}{88.21--92.32}
\newcommand{\NoneCommon}{906}
\newcommand{\NoneFone}{0.873}
\newcommand{\NoneKappa}{0.872}
\newcommand{\NoneMacro}{83.04}
\newcommand{\NoneMcc}{0.879}
\newcommand{\NoneRare}{17}
\newcommand{\NoneRareMacro}{77.0}
\newcommand{\NoneRci}{0.788}
\newcommand{\NoneSeedMean}{90.69}
\newcommand{\NoneSeedSd}{1.86}
\newcommand{\PairedCi}{2.47--5.76}
\newcommand{\PairedGain}{4.11}
\newcommand{\SeedDiff}{-0.35}
\newcommand{\SeedDiffSd}{2.77}
\newcommand{\TimingSeconds}{207.58}
% END VERIFIED NUMBERS

\title{Few-shot OCT Classification of Rare Retinal Diseases\\
  \large Current results, problems requiring discussion, and research plan}
\author{Lam\\\small Project working report}
\date{Evidence checkpoint: 27 October 2026}

\begin{document}
\maketitle

\textbf{Assessment.} I have implemented the main independent-test pipeline,
but I have not established the paper's central ordering of GAN augmentation
above basic augmentation across training seeds. My next priority is to complete
the benchmark re-evaluation and the agreed C0/C1 experiments. I then propose
extending the work with a focused comparison of accuracy and resource use
against a retinal foundation model.

In this report, I evaluate all five ideas in the research-directions artifact
inspected on 27 October 2026. I treat its novelty and feasibility ratings as
proposals, not measured findings. I draw the numerical results below from
saved artifacts and record the new descriptive calculations in
\file{docs/research/report_evidence.json}. I retain input SHA-256 hashes in that
record and provide \file{docs/research/build_report_evidence.py} to reproduce
it. I list the evidence IDs in the appendix and distinguish planned sample
sizes, decision thresholds and extrapolated runtimes from measurements.
\ev{E00}

\section{Guide to the work packages}

I use the labels C0--C5 for follow-up work on the Colab track, separately
from Stages 0--4 of the original data-to-classifier pipeline. The numbers
identify tasks, not a compulsory execution order. In
Table~\ref{tab:work-packages}, I explain each task's question, activity and
dependency. I use ``code ready'' to mean that an implementation is available,
not that the experiment has produced results.

\begin{longtable}{>{\raggedright\arraybackslash}p{0.06\linewidth}
                  >{\raggedright\arraybackslash}p{0.28\linewidth}
                  >{\raggedright\arraybackslash}p{0.33\linewidth}
                  >{\raggedright\arraybackslash}p{0.22\linewidth}}
\caption{What C0--C5 are for. Status is the saved project checkpoint, not
a claim about unsynchronized Drive runs. Sources: the replication tracker
and experiment plans. \ev{E08,E09,E10,E11}}
\label{tab:work-packages}\\
\toprule
Code & Question I address & Work I plan & Status and dependency\\
\midrule
\endfirsthead
\toprule
Code & Question I address & Work I plan & Status and dependency\\
\midrule
\endhead
\bottomrule\endfoot
C0 &
Does performance depend on file-format or image-framing cues? &
I plan to compare current images, uniform grayscale PNGs, and grayscale
images with square padding and resizing. I will extract features for each
input arm, check remaining separation, and compare runs if the gate permits. &
Code ready; sample diagnostics saved. I still need the full control before
choosing my inputs for C1.\\
C1 &
Does GAN augmentation reliably improve accuracy over basic augmentation
across classifier training seeds? &
I plan paired seed runs of \texttt{none}, \texttt{basic} and \texttt{gan},
with fixed250 final heads as primary and es20 as sensitivity. I will apply
the predeclared, single variance-based sample-size update. &
Protocol and code ready; no saved seed-pair outcomes. I will select inputs
after C0 and retain the original v1 GAN pool.\\
C2 &
How many synthetic GAN images are useful? &
I propose varying additional accepted GAN images per rare class across
0/500/1,000/2,000/3,000, with the basic pool fixed, and comparing across
seeds to address the paper's augmentation-amount experiment. &
Planned. I would run this after C5 supplies the full v2 pools, separately
from C1.\\
C3 &
Does reducing the dominance of common classes change performance? &
I propose under-sampling common-class training images to vary class imbalance
and compare classifier outcomes, addressing the paper's class-balance
experiment. &
Planned; I need a Stage 3 under-sampling option and will keep the test set fixed.\\
C4 &
Do conclusions persist across different training/evaluation folds? &
I propose fivefold cross-validation, rebuilding basic augmentation and
training CycleGANs within each training fold. I would exclude its evaluation
fold from all fitting and selection. &
Planned; I need more compute. I cannot reuse the original GANs across folds
without leaking information from held-out images.\\
C5 &
How can C2 compare equal, complete GAN-image budgets across rare classes? &
I plan to use existing GAN checkpoints to generate fresh candidates for
MH, MacTel, Stargardt and RP, apply the original automatic QC, and build
a separate v2 pool with 3,000 accepted images per class. &
Code ready; GPU execution pending. I will complete C5 before C2 and retain
C1's v1 pool.\\
\end{longtable}

\textbf{Execution sequence.} I plan to complete C0 and review its diagnostics
before selecting C1's input arm and running its fixed protocol. Separately,
I will complete C5 before C2. I regard C3 and C4 as later controls with their
own designs and compute requirements. For C4, I would respect known patient
and source groups; missing rare-patient identities remain a limitation.

\textbf{Related labels.} L7 is the local audit of the unexpectedly strong
real-only baseline; L8 is the RCI definition/numerical audit; L9 is the
patient-cluster uncertainty analysis. The labels v1 and v2 identify the
original and topped-up GAN pools. Fixed250 means 250 classifier epochs with
the final head retained; es20 means early stopping with patience 20 and the
best validation checkpoint retained. These are protocol names, not observed
performance improvements. \ev{E08,E10,E11}

\section{Current results}

\subsection{What I have reproduced}

I reproduce the independent-test experiment of Yoo, Choi and Kim
\cite{yoo2021}: NORMAL, CNV, DME and DRUSEN from OCT2017, together with CSC,
macular hole (MH), macular telangiectasia (MacTel), Stargardt disease and
retinitis pigmentosa (RP). Optical coherence tomography is abbreviated OCT.
The committed manifest contains 74,082 training, 9,526 validation and 1,022 test
images. The test set contains 250 images per common class and 5/5/4/4/4 images
for CSC/MH/MacTel/Stargardt/RP. Real rare-class training counts are
27/26/25/19/27, totaling 124. \ev{E01}

I use frozen ImageNet Inception-v3 features and train a softmax head. Throughout
this report, I use \texttt{none} for real rare training images, \texttt{basic}
for the existing basic-augmentation pool, and \texttt{gan} for basic
augmentation plus the accepted CycleGAN pool. These are my actual saved
training pools. For the augmented variants, I do not separately append the
original rare images. I use augmented derivatives for rare validation in
those variants; \texttt{none} has no rare validation images. \ev{E06}

My saved work covers Stages 0--4, the original GAN training/generation/QC,
seed-42 classifiers, the five-seed check, and ROC/t-SNE analysis data. I have
not replicated the entire paper: the augmentation-amount experiment,
under-sampling experiment and fold-specific fivefold cross-validation remain
incomplete. I implemented a class-mean baseline, not the paper's trained
prototypical network. I do not count C0/C1/C5 code readiness as experimental
completion. \ev{E08,E09,E10,E11}

\subsection{Training-seed variation changes the interpretation}

\begin{table}[htbp]
\centering
\caption{Saved accuracy across training seeds 1--5. SD is the sample SD
across five runs, not a confidence interval. All runs share the test set and
the saved augmentation/GAN realization.}
\label{tab:seeds}
\begin{tabular}{lrrl}
\toprule
Variant & Mean accuracy (\%) & SD (\pp) & Range (\%)\\
\midrule
\texttt{none} & \NoneSeedMean & \NoneSeedSd & 87.87--92.86\\
\texttt{basic} & \BasicSeedMean & \BasicSeedSd & 90.90--95.79\\
\texttt{gan} & \GanSeedMean & \GanSeedSd & 91.49--94.13\\
\bottomrule
\end{tabular}
\end{table}

The paired GAN-minus-basic accuracy differences, in seed order, were
$+2.94,-4.30,+0.10,+1.18,-1.66$ \pp. Their mean was
$\SeedDiff$ \pp{} with sample SD $\SeedDiffSd$ \pp;
GAN was ahead in three of five pairs. Basic and GAN had descriptively higher
means than \texttt{none}, by 2.246 and 1.898 \pp, respectively.
From these five pairs, I cannot establish GAN superiority or equivalence
between GAN and basic augmentation. Historical seed accuracies were saved to
two decimal places, and I retain that available precision. I do not have
committed per-image predictions for these five seeds, so I cannot estimate
their classwise intervals, calibration, or MCC variability here. \ev{E03}

\subsection{The seed-42 run of record}

\begin{table}[htbp]
\centering\small
\caption{The saved seed-42 classifiers. Accuracy intervals are the saved
patient-cluster percentile intervals. The paper column is its Table 3,
not a result from my reproduction. MCC denotes Matthews correlation coefficient.}
\label{tab:seed42}
\begin{tabularx}{\linewidth}{Yrrrr}
\toprule
Variant & Accuracy (\%) & 95\% CI (\%) & MCC & Paper accuracy (\%)\\
\midrule
\texttt{none} & \NoneAcc & \NoneCi & \NoneMcc & 88.4\\
\texttt{basic} & \BasicAcc & \BasicCi & \BasicMcc & 91.3\\
\texttt{gan} & \GanAcc & \GanCi & \GanMcc & 92.1\\
\bottomrule
\end{tabularx}
\end{table}

Using the evidence-audit script, I recompute the saved confusion matrices from
all three prediction CSVs, check test paths and labels against the manifest,
and independently recalculate accuracy, Cohen's $\kappa$, MCC and relative
classifier information (RCI). For \texttt{none}/\texttt{basic}/\texttt{gan}, $\kappa$ is
\NoneKappa/\BasicKappa/\GanKappa, and RCI is
\NoneRci/\BasicRci/\GanRci. The corresponding nine-class macro-F1 values,
derived from the saved confusion matrices, are
\NoneFone/\BasicFone/\GanFone. Balanced accuracy, defined here as unweighted
mean class recall, is \NoneMacro/\BasicMacro/\GanMacro\%.
I log these additional descriptive summaries in E00. \ev{E00,E02}

At seed 42, GAN exceeded basic by $+\PairedGain$ \pp,
with a saved paired patient-cluster 95\% CI of \PairedCi{} \pp.
The bootstrap used 5,000 draws with RNG seed 0, resampling 633 common-class
patient clusters and 22 assumed rare-image singleton clusters.
This interval conditions on the fitted heads and existing image pools. It
does not describe training-seed uncertainty. Rare patient identities are
unknown; shared rare patients could make these intervals too narrow. The
pairwise results are exploratory and not adjusted for multiple comparisons.
\ev{E04}

\subsection{Rare-class performance and where the gain comes from}

\begin{table}[htbp]
\centering\small
\caption{Rare-class recall at seed 42. Parentheses in the first column give
the number of test images. Macro recall weights the five diseases equally.}
\label{tab:rare}
\begin{tabular}{lrrr}
\toprule
Class & \texttt{none} & \texttt{basic} & \texttt{gan}\\
\midrule
CSC (5) & 0.800 & 0.800 & 0.800\\
MH (5) & 0.800 & 1.000 & 1.000\\
MacTel (4) & 0.750 & 0.500 & 0.750\\
Stargardt (4) & 0.750 & 1.000 & 1.000\\
RP (4) & 0.750 & 0.750 & 0.750\\
\midrule
Correct / 22 & \NoneRare & \BasicRare & \GanRare\\
Rare macro recall (\%) & \NoneRareMacro & \BasicRareMacro & \GanRareMacro\\
\bottomrule
\end{tabular}
\end{table}

One rare test image changes its disease recall by 0.20 or 0.25. GAN and basic
therefore differ by only one correctly classified rare image at seed 42.
Of GAN's 42 additional correct predictions over basic, 41 are common-class
images: common correct counts are \NoneCommon/\BasicCommon/\GanCommon.
DRUSEN correct counts change from 180 to 213 between basic and GAN, while
DME changes from 213 to 229. I therefore do not interpret the headline
overall gain as evidence of a broad rare-disease rescue. \ev{E02}

The paper's real-only recalls imply 3/22 rare images correct, compared with
17/22 in my reproduction. Its initial rare training counts total 97, compared
with my 124. My recorded training ran 26/27/22 epochs and selected minimum-validation-loss
checkpoints at epochs 6/7/2. The paper states 250 classifier epochs and a frozen
backbone, but does not establish my assumed final-head selection rule.
I cannot isolate the individual effects of these comparability differences.
The L7 audit reports no exact rare-image duplicates, but its original untracked
pixel-audit files are unavailable in this checkout. I retain that historical
finding without presenting it as a newly repeated duplicate audit.
\ev{E06,E08}\cite{yoo2021}

\subsection{Synthetic-image QC and visual observations}

\begin{table}[htbp]
\centering\small
\caption{Original automatic feature-distance QC, independently checked against
the saved per-image QC CSVs. Each disease has 4,000 candidates. Passing and
retained counts differ for CSC because retention is capped at 3,000.}
\label{tab:qc}
\begin{tabular}{lrrr}
\toprule
Disease & Passed & Retained & Pass rate (\%)\\
\midrule
CSC & 3,370 & 3,000 & 84.250\\
MH & 2,610 & 2,610 & 65.250\\
MacTel & 1,248 & 1,248 & 31.200\\
Stargardt & 2,199 & 2,199 & 54.975\\
RP & 832 & 832 & 20.800\\
\bottomrule
\end{tabular}
\end{table}

Automatic acceptance means that an image is closer to real rare-class images
than to NORMAL in frozen Inception-v3 feature space. It does not establish
clinical plausibility. Four classes lack the paper's target of 3,000 accepted
GAN images. I will preserve this original pool (v1) in C1 and use the prepared
C5 top-up to create a separate v2 pool for C2 only. I have dropped expert
GAN-image ratings from the replication because the required ophthalmologist
raters are unavailable. \ev{E05,E11}

In my saved, unblinded, non-expert CAM review, I judged the strongest activation
to lie on the abnormality in 21/22 GAN maps and 14/22 real-only maps; my off-retina
responses were 0/22 and 7/22. I treat these as my observations about saliency,
not lesion-localization ground truth or proof of pathology-based reasoning.
I retain the review CSV but have not recreated the underlying untracked
review assets for this report. \ev{E07}

\section{Problems requiring discussion}

\subsection{Patient overlap limits the central performance claim}

The manifest identifies 633 common-class test patients. Of these, 566 patients
accounting for 914/1,000 test images also occur in train or validation;
506 patients accounting for 804 test images occur in the fitting split alone.
I independently verify these counts with the evidence script. Patient-cluster
resampling addresses dependence within the test set; it cannot remove
train--test patient overlap. I therefore interpret my current accuracy as
performance on the committed benchmark, not as an estimate of generalization
to unseen patients.
\ev{E01,E04}

\textbf{Decision needed.} I will preserve the existing manifest as my replication
record. I would design any patient-disjoint benchmark as a separately versioned
experiment and fix patient/source identities and duplicate-group assignments
before model fitting. I would maintain that separation in fitting, validation,
augmentation, generator training and QC references. Rare filenames do not
supply verified patient IDs, so I cannot certify the present rare split as
patient-disjoint.

\subsection{Dataset source and disease are entangled}

The saved diagnostic reports that native image width/height alone identifies
16/22 rare test labels, and that a rare-versus-common feature axis has test AUC
0.9955. These diagnostics cannot distinguish disease morphology from
acquisition or source effects. The C0 delivery report documents a bounded
sample check: framing reduced the recomputed-axis AUC to 0.7125, but the fixed
original axis retained AUC 0.9250. Those sample results are recorded in the
delivery document; its referenced untracked \file{cue_check.json} is absent
here. They are not full-cohort C0 results. \ev{E08,E09}

\textbf{Decision needed.} I need to complete C0's full feature diagnostic before
selecting the C1 input arm. The predeclared gate stops before head training if framing's
full recomputed-axis AUC is at least 0.90. Passing this gate does not establish
removal of source cues. I will also review the fixed-axis diagnostic and
describe C0 as a preprocessing sensitivity experiment unless I obtain stronger
source-matched evidence.

\subsection{Seed uncertainty, protocol differences and metric comparability}

C1 addresses classifier randomness and checkpoint schedule on the fixed
benchmark. It cannot resolve unknown patient identities, historical feature
provenance, generator-seed variation or the small test sample. Its seeds also
change the augmented validation draw, so their variation includes more than
head initialization. Matching cache path lists verifies identities and order,
but not the historical backbone weights, input bytes or extraction environment.
\ev{E06,E08,E10}

I define RCI in this repository as $I(Y;\widehat{Y})/H(Y)$ and independently
reproduce its saved values with the report audit. The saved L8 analysis gives a
conditional upper bound of 0.970242 for the paper's GAN Table 4 recalls under
this definition, below the reported 0.983. The bound is a saved analytical
result; I have not rerun that optimization for this report. The paper's actual RCI
implementation remains unknown. Its Table 4 also implies prototypical-network
accuracy 86.1\%, whereas Table 3 reports 80.8\%. These inconsistencies warrant
clarification, not an invented replacement formula. I have kept author
correspondence on hold. \ev{E08}\cite{yoo2021}

The saved K=3 Euclidean kNN accuracy is 86.40\%, versus the paper's 91.2\%.
Changing to cosine K=1 reaches 91.10\% in the saved exploratory sweep, but that
changes the specified method. The exact feature layer and postprocessing in
the original baseline remain unresolved. \ev{E08}

\subsection{The extension needs a different evidence standard}

Overall accuracy is dominated by common classes. For a future rare-disease
study, I propose predeclaring rare macro recall as the principal performance
outcome and reporting disease-level denominators, common-class performance,
macro-F1, balanced accuracy and MCC. For C1, I retain the already agreed
primary outcome of overall accuracy. To assess calibration, I would need saved
probabilities, a fixed calibration procedure and validation-only tuning.
I have no measured expected-calibration-error or FID result to report here.

Before making clinical or resource claims, I need additional records.
GPU-hours alone do not describe CPU extraction, generator training, rejected
images or expert time. I cannot treat automatic QC pass rates as expert
acceptance rates. For a dollar-cost comparison, I would need the actual
hardware, measured usage, billing basis/date and an explicit treatment of
pretrained-model cost. I currently have no measured comparative dollar cost,
energy saving, expert-review time or foundation-model accuracy to report.

\section{Evaluation of the five proposed open gaps}

\subsection{Idea 2: rigorous benchmark re-evaluation --- first priority}

\textbf{Value.} I prioritize this direction because my findings already show
patient overlap, limited rare test support, source-sensitive features and
unstable seed ordering. I aim to provide a reproducible account of which
augmentation conclusions survive explicit controls. Because confidence
intervals and exact-duplicate checks are partly completed, I do not present
them as entirely new work.

\textbf{Remaining gap.} Similarity screening is not verified patient identity.
Near-duplicate candidates need source tracing and adjudication. C0 cannot
certify clinical generalization, and C1 cannot remove leakage. I propose
completing the remaining controls before adding a separately named
patient-disjoint evaluation, if suitable metadata exist. My aim is an auditable
validity boundary and a stable comparison, even if GAN does not outperform
basic augmentation.

\subsection{Idea 1: cost versus rare-class performance --- second priority}

\textbf{Value.} I regard this as the strongest extension after re-evaluation
if I narrow the question to the marginal resource cost of useful rare-class
performance. My minimal proposed design crosses frozen representation
(ImageNet Inception-v3 versus RETFound-OCT) with training pool
(basic versus basic plus GAN), retaining the real-only baseline as context.
This would separate the augmentation contrast within each backbone. I would
hold patient split, support identities, synthetic counts, tuning budget and
evaluation rules fixed, and pair runs by seed and support selection.

\textbf{Novelty boundary.} A newer generator by itself is a weak contribution.
Tai et al. already combine U-GAT-IT, balancing and attention; Iwanicka and Lu
study class-aware diffusion with adversarial fine-tuning; EyeDiff already
examines synthetic augmentation of retinal foundation models across multiple
datasets. I would base my contribution on the controlled resource/accuracy
comparison and its failure conditions. Before claiming to be the first such
study, I would need a dedicated, updated literature search.
\cite{tai2025,iwanicka,eyediff}

\textbf{Feasibility corrections.} The official Nature RETFound-OCT weights are
access-gated and carry a non-commercial license. Public availability does not
demonstrate low adaptation cost. More seriously, the RETFound paper explicitly
includes Kermany OCT data in self-supervised pretraining. Exact overlap with
my test files is unverified. I therefore need a pretraining-exposure audit
before treating a common-class result on this benchmark as an unseen-image
comparison. I would report marginal adaptation cost separately from the
external pretraining resource. Reported label efficiency in other tasks
does not guarantee the same behavior for these rare diseases.
\cite{retfound,retfoundweights}

The proposed 20-shot endpoint is not available for every class in the current
fitting split: Stargardt has 19 real training images. A support curve must use
feasible distinct real examples and preserve the holdout; augmented copies do
not create additional real shots. Defining shots as patients would require
patient identities that are currently missing. Diffusion and fine-tuning arms
would be later extensions, after I establish whether the frozen-feature
comparison is useful and decide on the scope and compute budget. \ev{E01}

\subsection{Idea 5: local clinical external validation --- start feasibility early}

\textbf{Value.} I see this as the most direct route to learning whether my
model works outside the web-image benchmark. It could improve the scientific value
of either re-evaluation or the cost comparison, but it is not currently a
ready experiment. I do not yet have a verified hospital dataset, recruitment
agreement, patient counts or expert labels in my project records.

\textbf{Requirements.} I would need a clinical partner, applicable approvals,
patient/eye/device metadata, label adjudication and a test-only use plan before
claiming an external result. I would fix models and thresholds before opening
that test, and plan independent patient counts per disease from the desired
precision instead of adopting the artifact's illustrative scan-count range.
I have not verified the suggestion that this would be the first clinic-sourced test;
the existing EyeDiff work already includes external validation in related
settings. A small convenience cohort would be a feasibility study, not
clinical validation at deployment scale. \cite{eyediff}

\subsection{Idea 3: open-set detection --- useful, but defer the main experiment}

\textbf{Value.} I consider unseen-disease detection a meaningful question,
but my present nine-class softmax outputs provide no open-set performance
estimate. Low confidence, prototype distance or an energy score are candidate
rules, not validated unknown-disease detectors.

\textbf{Discriminating design.} I would hold a disease out of all supervised
fitting, augmentation and selection stages, and choose rejection thresholds
using separate validation identities and unknown categories. I would evaluate
known-class recall alongside unknown-detection sensitivity and false-positive
rate at the predeclared threshold. I would also keep known and unknown
acquisition sources comparable; otherwise the detector could merely recognize
the source. My small rare holdout and unknown patient identities limit this
experiment, so I would treat it as a separate follow-up once the benchmark is
stable.

\subsection{Idea 4: one conditional generator --- lowest immediate priority}

\textbf{Value.} A shared generator might reduce incremental cost as diseases
are added, but I have not measured this. Class-conditioned diffusion
and multimodal retinal generation already exist, which narrows the novelty
of replacing five separate CycleGANs with one model. \cite{iwanicka,eyediff}

\textbf{Decision rule.} I would pursue this only if my cost benchmark identifies
per-disease generator training or review as the limiting expense. I would
compare matched real-image budgets, compute, accepted synthetic counts and
downstream rare recall as the disease set grows, and audit class-specific
failures and memorization. A shared model can also introduce negative transfer;
fewer model files alone do not establish lower total cost or better diagnosis.

\subsection{OCTDL is a complementary task, not a replacement for missing labels}

The OCTDL release described by Kulyabin et al. contains 2,064 images from
821 patients and explicit patient IDs. Its RAO group has 22 images from
11 patients. Its categories are AMD, DME, ERM, NORMAL, RAO, RVO and VID;
it does not supply the same five rare-disease labels as my present experiment.
I consider it promising for a separately specified patient-disjoint benchmark,
but I cannot use it to directly validate all nine existing outputs. \cite{octdl}

AMD does not map one-to-one onto CNV or DRUSEN, and the OCTDL description
includes macular-hole pathology within its condition metadata. Hence a
category absent by name is not automatically a clinically unseen pathology.
I would define label mappings with expert input before pooling datasets or
declaring open-set classes. More scans also do not guarantee an adequately sized rare
test cohort: independent patient counts remain the relevant constraint.
\cite{octdl}

\section{Planning proposals for the upcoming weeks}

\subsection{My experimental sequence}

I propose the following work windows; they are not completed runs or a guarantee
of Colab availability. I summarize task definitions and dependencies in
Table~\ref{tab:work-packages}. I will review C0 before selecting C1's inputs,
retain the original v1 GAN pool for C1, and reserve C5's top-up for C2.
I need the code on GitHub before running it in Colab. I will make separate
decisions about publishing code, new model dependencies and new dataset/split
protocols. \ev{E09,E10,E11}

\begin{longtable}{>{\raggedright\arraybackslash}p{0.16\linewidth}
                  >{\raggedright\arraybackslash}p{0.43\linewidth}
                  >{\raggedright\arraybackslash}p{0.32\linewidth}}
\caption{My four-week proposal. Each later window depends on its stated
prerequisites; I will carry incomplete work forward while retaining the analysis rule.}
\label{tab:weeks}\\
\toprule
Window & My proposed work & Planned output / gate\\
\midrule
\endfirsthead
\toprule Window & My proposed work & Planned output / gate\\\midrule
\endhead
\bottomrule\endfoot
Week 1\\28 Oct--3 Nov &
I plan to freeze the evidence inventory, complete the full C0 input/feature
diagnostics, review source and duplicate candidates, and confirm Colab/Drive
readiness. I could also begin clinical-partner feasibility discussions and
review RETFound access. &
C0 \file{full_axis.csv}, \file{gate.json}, inventories and source-audit notes.
If the framing gate fails, I will stop C0 heads and decide on the next step.
I will record which patient/source claims remain unverifiable.\\
Week 2\\4--10 Nov &
If C0 permits, I will run its paired heads and summarize all contrasts.
I will select C1's input arm before seeing C1 outcomes, then start the agreed
first 20 seed blocks under fixed250 and es20 across all three variants. &
C0 decision record; C1 identities, commands, timings and complete run folders.
I will finish all first-stage protocol/variant runs before the single
variance recalculation and label partial summaries as incomplete.\\
Week 3\\11--17 Nov &
I plan to finish the first C1 stage if necessary, lock the variance-based
top-up target, and run fixed250 seeds 21 through that target. I will also draft
the narrower foundation-model comparison, support budgets and cost ledger
for review. &
Immutable \file{topup.json} and run schedule. If the target exceeds the
predeclared cap, I will report infeasibility. I will begin new model work only
under a separately accepted protocol.\\
Week 4\\18--24 Nov &
I will summarize C1 only if the locked schedule is complete; otherwise I will
document progress and remaining compute. If the extension protocol and access
are ready, I will begin its isolated feature/timing pilot; otherwise I will
finalize the design. I will place C5 before any later C2 dose experiment. &
Paired contrasts with uncertainty and the zero/1-pp interpretation; an updated
evidence ledger. Any pilot supplies measured timing and feasibility, with no
claim of completed foundation-model or clinical validation.\\
\end{longtable}

\subsection{C1 targets are protocol choices, not results}

My fixed C1 plan makes final heads after 250 epochs the primary analysis
and early-stopping patience 20 the sensitivity analysis on seeds 1--20.
The primary contrast is mean GAN-minus-basic overall accuracy; the meaningful
target is 1 \pp. I will complete both protocols and all variants for the first
20 seeds, then perform one SD-only paired-$t$ power recalculation at two-sided
$\alpha=0.05$ and target power 0.80. The total target has a floor of 63 paired
seeds and a feasibility cap of 243. These numbers are predeclared planning
rules. There are no saved C1 seed-pair outcomes at this checkpoint. \ev{E10}

If the estimated target exceeds 243, I will report the study as beyond the
declared budget. Otherwise I will complete the locked fixed250 schedule.
The final paired-$t$ interval is nominal under this variance-adaptive sample
size. I will follow the plan without significance-based stopping or repeated
sample-size adjustment, and report the interval relative to both zero and
1 \pp; detecting a nonzero effect is not the same as
establishing an effect of at least 1 \pp. \ev{E10}

\subsection{My runtime budget from measured timing}

I base my runtime budget on the saved local CPU timing of 25 completed
cached-head epochs in \TimingSeconds{} seconds. Its linear extrapolation is
\HeadMinutes{} minutes for 250 epochs, including repeated startup overhead.
My existing plan projects approximately 37--42 serial CPU hours for the first
C1 stage and 112--116 hours for a 63-pair primary study including es20
sensitivity. Those totals assume 20--50 epochs for the es20 heads; I treat them
as forecasts, not measured Colab durations. I have not verified Colab CPU/GPU
speed or session limits.
\ev{E10}

The original Stage 2 log separately records 104/116/115/126/133 seconds to
generate 4,000 candidates for CSC/MH/MacTel/Stargardt/RP on the logged A100
runtime. Their sum is 594 seconds for generation only. It excludes GAN
training, QC, feature extraction, transfer and human review, so I cannot use
it as the total cost of generative augmentation. I also cannot assume that
the logged training hardware and software match a future foundation-model run.
\ev{E05}

\subsection{Records I will retain from each experiment}

I will retain one immutable run identity: code revision, environment versions,
manifest and input/cache hashes, preprocessing, backbone/checkpoint hash,
pool version, real/synthetic counts, support selection, seed and complete
command. I will save all test probabilities and labels, selected/best/final
heads, history, timing, log and a completion marker, and distinguish
patient-level uncertainty from training-seed uncertainty.

For the cost comparison, I will additionally log hardware, CPU/GPU elapsed
usage, feature-extraction and generator time, candidate/retained counts,
storage and expert time if actually observed. I will record the billing
source and date for any USD conversion and retain missing measurements as
missing. I will use a new versioned output namespace for the extension and
preserve canonical Stage 3/4 and \file{results/seed_check.csv}. \ev{E10,E11}

\textbf{My decisions for the next meeting.} I need to select C1's input arm
after C0, confirm the multi-session compute budget, and decide whether my
next deliverable covers benchmark validity alone or also includes the narrow
cost comparison. I also need to establish whether patient-identified external
data and a clinician collaborator are realistically available. I have not
initiated author contact about the unresolved paper details as part of
preparing this report.

\appendix
\section{Evidence locations and reproduction}

I give all paths below relative to the repository root. In preparing this
report, I add descriptive summaries and interpretation while preserving the
method, manifest, accepted images and published evaluation CSVs. I take C0,
C1 and C5 execution status from the current plans and available saved
artifacts; I cannot verify the content of unsynchronized private Drive folders.

\begin{longtable}{>{\raggedright\arraybackslash}p{0.09\linewidth}
                  >{\raggedright\arraybackslash}p{0.84\linewidth}}
\toprule ID & Saved evidence and its scope\\\midrule
\endfirsthead
\toprule ID & Saved evidence and its scope\\\midrule
\endhead
\bottomrule\endfoot
E00 & \file{docs/research/report_evidence.json}; generated/checked by
\file{docs/research/build_report_evidence.py}. Contains source hashes,
unrounded descriptive calculations and numerical LaTeX values. The companion
\file{docs/research/report_sources.json} records the read proposal and primary
source locations.\\
E01 & \file{data/manifest.csv}: splits, real training/test counts, image
identities and common-class patient overlap.\\
E02 & \file{results/stage3_<variant>/test_predictions.csv};
\file{results/stage4/metrics.csv}, \file{recall.csv},
\file{confusion_<variant>.csv}: seed-42 outcomes and independently recomputed
descriptive metrics.\\
E03 & \file{results/seed_check.csv};
\file{results/analysis/paired_seeds.csv}: historical seeds 1--5 and paired
differences. These are distinct from the planned C1 runs.\\
E04 & \file{results/analysis/cluster_seed42.csv},
\file{paired_cluster_seed42.csv}, \file{cluster_summary.csv}:
saved patient-cluster intervals and overlap summary. I do not regenerate
bootstrap draws in the report audit.\\
E05 & \file{results/stage2/qc_<disease>.csv};
\file{results/analysis/qc_summary.csv}; \file{results/stage2/stage2.log}:
automatic QC and measured generation timings.\\
E06 & \file{results/stage3_<variant>/data_counts.csv}, \file{history.csv};
\file{stage3_classifier.py}: actual pools, validation composition,
checkpoint selection and classifier implementation.\\
E07 & \file{results/analysis/cam_review_seed42_Lam.csv}: all 44 of my CAM
responses. I do not treat these as expert ratings.\\
E08 & \file{docs/research/L7_BASELINE_AUDIT.md},
\file{L8_RCI_AUDIT.md}, \file{PAPER_DISCREPANCIES.md},
\file{REPLICATION.md}; \file{results/analysis/rare_recall.csv},
\file{rci_paper_bounds.csv}, \file{knn_sweep.csv}:
bounded historical audits and unresolved comparisons. I do not treat their
absent local scratch files as independently re-executed evidence.\\
E09 & \file{docs/research/C0_DESIGN.md}, \file{C0_REPORT.md};
\file{c0_experiment.py}; notebook Part C0: sample diagnostics, gate and pending
full control. The sample JSON referenced by C0 is absent locally.\\
E10 & \file{docs/research/C1_EXPERIMENT_PLAN.md},
\file{C1_HEAD_TIMING.json}; \file{c1_experiment.py}; notebook Part 4:
fixed250/es20 protocol, one variance recalculation, output protection and CPU
timing. Forecasts are labelled separately from measurements.\\
E11 & \file{COLAB.md}, C5 section; \file{c5_topup.py};
\file{docs/research/CONTEXT.md}: agreed v1/v2 separation and pending GPU top-up.\\
\end{longtable}

I provide the following command to reproduce the training-free evidence check
from this Windows worktree:
\begin{quote}\small\ttfamily
\detokenize{D:/Few-shot/.venv/Scripts/python.exe -X utf8}\par
\detokenize{docs/research/build_report_evidence.py check}
\end{quote}
The two displayed lines form one command. The \texttt{write} mode deliberately
refreshes only the report evidence record and the report's numerical macros.
To compile the self-contained source, I would run \texttt{pdflatex -output-directory=docs/report docs/report/Report.tex}
twice in a TeX installation with the packages listed in the preamble.
I have checked the source structure, but final compilation and page layout
remain unverified without a TeX engine.

\textbf{Terminology.} I use ``GAN'' for the original CycleGAN pool unless I
explicitly name a new version. Accuracy and balanced accuracy are percentages;
differences are \pp; recall, macro-F1, $\kappa$, MCC, RCI and AUC are fractions
unless a table labels them as percentages. A classifier seed is a training
replicate, not a new patient cohort. A QC pass is an automatic decision,
not a clinician's acceptance.

\textbf{Preparation note.} I used agent assistance to organize the evidence,
independently check saved-result arithmetic and draft this report. I also used
Scientific Agent Skills to guide the appraisal and planning workflow
\cite{skills}; these are procedural aids, not evidence for biological or
clinical conclusions. I remain responsible for the research decisions and
final interpretation.

\begin{thebibliography}{9}
\bibitem{yoo2021} Yoo TK, Choi JY, Kim HK. Feasibility study to improve deep
learning in OCT diagnosis of rare retinal diseases with few-shot classification.
\emph{Medical \& Biological Engineering \& Computing}. 2021;59:401--415.
\href{https://doi.org/10.1007/s11517-021-02321-1}{doi:10.1007/s11517-021-02321-1}.
Methods and Tables 3--4 checked in the full Europe PMC XML.

\bibitem{tai2025} Tai CY et al. Improvement Strategies for Few-Shot Learning in
OCT Image Classification of Rare Retinal Diseases. 2025. Preprint.
\href{https://arxiv.org/abs/2505.20149}{arXiv:2505.20149}.
Abstract and full-text methods checked.

\bibitem{iwanicka} Iwanicka D, Lu P. Diffusion with Adversarial Fine-Tuning for
Improving Rare Retinal Disease Diagnosis. MIUA 2025; proceedings volume
published 2026. LNCS 15918:237--250.
\href{https://doi.org/10.1007/978-3-031-98694-9_17}{doi:10.1007/978-3-031-98694-9\_17}.
\href{https://eprints.whiterose.ac.uk/id/eprint/228674/}{Author-repository abstract and metadata checked};
full experimental methods not audited here.

\bibitem{eyediff} Chen R et al. Boosting foundation models for rare eye disease
diagnosis via a multimodal text-to-image generative framework.
\emph{npj Digital Medicine}. 2026;9:371.
\href{https://doi.org/10.1038/s41746-026-02560-2}{doi:10.1038/s41746-026-02560-2}.
Publisher abstract and relevant Results sections checked; not an exhaustive
replication or cost audit.

\bibitem{retfound} Zhou Y et al. A foundation model for generalizable disease
detection from retinal images. \emph{Nature}. 2023;622:156--163.
\href{https://doi.org/10.1038/s41586-023-06555-x}{doi:10.1038/s41586-023-06555-x}.
Full-text pretraining dataset description and reference 34 checked in Europe PMC XML.

\bibitem{retfoundweights} RETFound authors.
\href{https://huggingface.co/YukunZhou/RETFound_mae_natureOCT}{Official Nature OCT checkpoint model card}
and \href{https://github.com/rmaphoh/RETFound}{official repository}.
Access gate and model-card license checked 27 October 2026; weights not downloaded.

\bibitem{octdl} Kulyabin M et al. OCTDL: Optical Coherence Tomography Dataset
for Image-Based Deep Learning Methods. \emph{Scientific Data}. 2024;11:365.
\href{https://doi.org/10.1038/s41597-024-03182-7}{doi:10.1038/s41597-024-03182-7}.
Full-text Data Records, label definitions and Table 2 checked in Europe PMC XML;
dataset files not downloaded or counted locally.

\bibitem{skills} Kassis T, Agarwal V, He Y, Patel D, Brueckner AM.
Scientific Agent Skills: A Library of Procedural Knowledge for Research Agents.
2026. \href{https://doi.org/10.48550/arXiv.2609.00065}{arXiv:2609.00065}.
Author/year metadata checked against the current arXiv record on 27 October 2026.
\end{thebibliography}
\end{document}
